\section{Experimental setup}
\label{sec:experiments}

\subsection{Models, data and grid}

We evaluated four convolutional backbones: ResNet-34 and ResNet-50
\citep{he2016deep}, and VGG-11 and VGG-19 \citep{simonyan2015very}. The pair
structure is deliberate --- two depths within each family --- because the two
families differ in exactly the properties we expected to matter. The ResNets are
heavily normalised (ResNet-50 carries 53 BatchNorm layers) and hold most of their
parameters in convolutions; both VGGs carry no BatchNorm at all and hold
approximately 85\% of their parameters in the classifier MLP. Section~\ref{sec:experiments:vgglr}
reports where that difference surfaced.

All experiments used CIFAR-10 \citep{krizhevsky2009cifar}. We split the 50{,}000
image training set 9:1 into \textbf{45{,}000 training} and \textbf{5{,}000
validation} images, and held out the \textbf{10{,}000 image test set} for final
evaluation only. At batch size 512 this gives 87 optimizer steps per epoch. All
reported test accuracies are over the full 10{,}000 image test set; no evaluation
in this paper drops a partial final batch.

The grid was 4 models $\times$ 3 pruning criteria $\times$ 4 sparsity levels
(0.95, 0.97, 0.99, 0.999) $\times$ 2 arms, plus 4 dense reference models: 100
cells in total.

\subsection{Protocol}

Table~\ref{tab:protocol} gives the full protocol. It is written as three columns
because there are three training phases, but the important reading of it is
horizontal: the I.P.\ and BaCP columns are identical on every row except the
last.

\begin{table}[t]
\centering
\small
\caption{Training protocol. The I.P.\ and BaCP arms are identical on every row
except \emph{objective}. Both arms received the same initialisation, the same
sparsity schedule, the same mask-update interval, the same epoch budget and the
same fine-tune; this matched budget is what makes the I.P.--BaCP difference
attributable to the objective.}
\label{tab:protocol}
\begin{tabular}{llll}
\toprule
 & Dense reference & I.P.\ baseline & BaCP \\
\midrule
Initialisation      & ImageNet-pretrained & dense checkpoint & dense checkpoint \\
Optimizer           & SGD, LR 0.01 & SGD, LR 0.01 & SGD, LR 0.1 (VGG: 0.05) \\
Epochs              & 100, patience 20 & 50 & 50 \\
Fine-tune           & --- & AdamW $10^{-4}$, 25 ep. & AdamW $10^{-4}$, 25 ep. \\
Mask interval $\Delta T$ & --- & 100 steps & 100 steps \\
Sparsity schedule   & --- & cubic, 80/20 ramp/hold & cubic, 80/20 ramp/hold \\
Batch size          & 512 & 512 & 512 \\
Precision           & bf16 autocast & bf16 autocast & bf16 autocast \\
Gradient clipping   & 10.0 global norm & 10.0 global norm & 10.0 global norm \\
Objective           & CE & CE & $\lambda$-weighted PrC+FiC+SnC+CE \\
\bottomrule
\end{tabular}
\end{table}

\paragraph{The matched-budget argument.}
This is the central methodological commitment of the paper, so we state it
without hedging. The I.P.\ baseline is not a separately-tuned competitor method.
It is the \emph{same run} as BaCP --- same pretrained initialisation, same dense
checkpoint, same pruning criterion, same cubic schedule, same $\Delta T$, same
prunable scope, same 50-epoch pruning budget, same 25-epoch AdamW fine-tune, same
data splits, same precision, same gradient clipping --- with the four-term
objective of Eq.~\eqref{eq:objective} replaced by plain cross-entropy. In
particular, the fine-tune is applied to \emph{both} arms. Without that, the
comparison degenerates into 75 epochs of BaCP against 50 of I.P., and a reader
cannot separate the objective from the extra training. With it, the objective is
the only free variable, and any difference between the arms is attributable to it
and to nothing else. \citet{blalock2020state} document how often pruning
comparisons fail precisely here; the design above is our response to that.

\paragraph{Model selection.}
BaCP's contrastive phase has no supervised validation signal to select on --- its
objective is dominated by terms that are not a task metric --- so the checkpoint
carried out of that phase was selected on \emph{training loss}. Both arms'
fine-tuning phases selected on \emph{validation accuracy}, computed on the held-out
5{,}000 image validation split. \textbf{No stage of either arm selected on the
test set.} The test set was evaluated once per completed cell.

\paragraph{Precision, clipping and hardware.}
All training ran under bfloat16 autocast with fp32 master weights. Gradients were
clipped at a global norm of 10.0 on \emph{both} arms; the value is SNIP-it's
published one \citep{verdenius2020snipit}, and applying it symmetrically matters
because an asymmetric stabiliser would itself be a difference between the arms.
Runs were executed one at a time on a single NVIDIA A100 (Azure
\texttt{Standard\_NC24ads\_A100\_v4}, 24 vCPU), with the data loaders given the
full CPU allocation.

\paragraph{Seeds.}
Run-to-run variation on this task and protocol was measured at approximately
\result{cifar10.noise.floor.lo}--\result{cifar10.noise.floor.hi} accuracy points.
Because that floor is not negligible relative to the effects we are looking for,
the headline cells were run with \textbf{3 seeds} and are reported as mean
$\pm$ standard deviation; the remaining cells of the grid are single-seed. We
mark which is which in every table. This is a compromise forced by compute, and
we count it as a limitation in Section~\ref{sec:discussion} rather than presenting
partial seed coverage as adequate \citep{bouthillier2021variance}.

\subsection{A protocol detail: the VGG learning rate}
\label{sec:experiments:vgglr}

The BaCP arm used SGD at LR 0.1 for the ResNets and 0.05 for both VGGs. This is
the one place where the two arms' hyperparameters are not literally identical
across all four models, so we report how the value was chosen rather than
presenting it as a given.

Neither VGG contains a BatchNorm layer, and roughly 85\% of each VGG's parameters
sit in its classifier MLP. At LR 0.1 the contrastive phase destabilised on these
models: gradient clipping prevented divergence, but the resulting sparse models
sat below their own dense reference. We ran a direct probe on
\texttt{vgg19}/\texttt{magnitude}/0.95 with everything else held identical and
only the learning rate varied. BaCP at LR 0.10 reached
\result{vgg19.probelr010.magnitude.0.95}; at LR 0.05 it reached
\result{vgg19.probelr005.magnitude.0.95}; the matched I.P.\ baseline reached
\result{vgg19.probeip.magnitude.0.95}. At 0.10 the objective was therefore worse
than the baseline it is compared against, and at 0.05 it was better. VGG-11
shares the architecture and exhibited the same symptom, so it took the same
value.

We report this because it is a real degree of freedom. Two things bound it. It
was fixed \emph{before} the reported sweep and applied uniformly to both VGGs; and
it was chosen against the I.P.\ baseline on a single cell, not swept per cell,
per criterion or per sparsity. The learning rate of the I.P.\ arm was not
correspondingly tuned, which is conservative in the direction that matters: the
baseline was left at the standard value while the method under test was made to
train stably.

\subsection{Relation to published comparators}
\label{sec:experiments:budget}

Section~\ref{sec:results} places our numbers next to published CIFAR-10 results
for the same architectures. Those comparisons need a stated caveat, because our
compute budget is substantially smaller. The GraNet CIFAR-10 recipe
\citep{liu2021granet}, which also supplies the GMP, SNIP, GraSP, SynFlow and RigL
rows we quote, trains for 160 epochs at batch 128 with a step learning-rate
schedule; ours is 50 pruning epochs plus 25 fine-tuning epochs at batch 512. In
optimizer updates that is close to an order of magnitude fewer. EAST's extreme
sparsity results \citep{li2025east} use 250 epochs at batch 128.

The consequence is specific and we accept it in advance: \textbf{our absolute
accuracies are expected to sit below the published ones, and no claim in this
paper is an absolute state-of-the-art claim.} What our protocol supports is the
\emph{controlled} quantity --- the difference between two arms that share a budget
--- and that quantity is unaffected by the budget being small, because both arms
are small in the same way. We quote the published numbers for calibration: if our
I.P.\ magnitude baseline lands within a plausible distance of published GMP under
a much shorter schedule, then the baseline BaCP is measured against is a real
baseline and not a strawman, and that is the only load the published numbers are
asked to carry here.
